\chapter{Metodologia}
\label{cap:metodologia}

O trabalho combina revisão de literatura, análise documental e desenvolvimento técnico. A escolha por essas etapas responde à natureza do problema: não basta entender a teoria do controle social, nem basta saber construir um Data Warehouse. É preciso fazer os dois conversarem em um caso concreto, com dados reais e suas inconsistências.

\section{Etapas}

\subsection*{Pesquisa Bibliográfica}

A revisão de literatura cobre quatro frentes que se cruzam neste trabalho: orçamento público brasileiro (instrumentos, classificações e estágios da despesa), transparência governamental e controle social, modelagem dimensional segundo Kimball e Ross, e \textit{civic tech} aplicada a dados públicos. A bibliografia orienta as decisões teóricas e fornece o vocabulário para interpretar os resultados. A relação entre literatura técnica e literatura sociopolítica é uma escolha deliberada e estabelece a tese central do trabalho: arquitetura de dados é objeto político.

\subsection*{Pesquisa Documental}

A análise documental tem como base o Portal da Transparência da Prefeitura de Juiz de Fora. Inclui o mapeamento das fontes disponíveis (planilhas, relatórios, APIs quando existem), o catálogo de campos por tipo de dado e a documentação das limitações encontradas.

\subsection*{Desenvolvimento Técnico do Data Warehouse}

O desenvolvimento segue a metodologia de Kimball em quatro passos: identificação dos processos de negócio (execução orçamentária de despesas e de receitas), definição da granularidade, identificação das dimensões relevantes (tempo, unidade administrativa, funcional programática, função e subfunção LOA, natureza da despesa, natureza da receita, fornecedor, fonte de recurso, empenho) e identificação dos fatos (valores empenhado, liquidado e pago para despesa; previsto e arrecadado para receita). A pilha técnica combina Dagster para orquestração, Python e \texttt{pandas} para extração, DuckDB como banco analítico embarcado, dbt para a modelagem em SQL e Streamlit para os \textit{dashboards}. As decisões técnicas detalhadas vivem no Capítulo~\ref{cap:dw}.

\subsection*{Estudo de Caso}

A aplicação do modelo a Juiz de Fora cumpre dois papéis. Primeiro, valida o modelo dimensional proposto: se o esquema responde às perguntas do mundo real, ele serve. Segundo, produz documentação dos desafios técnicos encontrados em um município específico, oferecendo referência para futuras adaptações em outras cidades. A escolha de Juiz de Fora não é casual e responde tanto à vinculação institucional do trabalho (UFJF) quanto à existência de portal municipal ativo com dados publicados em formatos minimamente analisáveis.